Retrieval-augmented generation separates a generator from an external evidence store \citep{lewis2020rag}; REALM studies retrieval-augmented language-model pretraining \citep{guu2020realm}. These are relevant foundations rather than interchangeable implementations. The present pilot uses a small local model and lexical retrieval over permitted derivative records, not the full architectures, data or training regimes of those papers. The intervention is access to source-addressable context under a controlled split.

\citet{ovadia2024finetuning} compare fine-tuning and retrieval for knowledge injection on their evaluated tasks. Their result motivates measuring both routes under stated conditions, not a universal ranking of RAG over every adaptation method. \citet{gekhman2024hallucinations} show why learning new factual associations during fine-tuning warrants separate analysis of unsupported outputs. Low-rank adaptation reduces the trainable parameter subset \citep{hu2022lora}; efficiency does not establish that learned responses are correct, faithful to source or robust to paraphrase.

RAG evaluation frameworks separate properties such as contextual relevance and answer faithfulness rather than assigning a single undifferentiated quality score \citep{es2024ragas,saadfalcon2024ares}. Their model-based evaluation components and calibration requirements cannot be inherited merely by computing token overlap. We retain raw outputs and retrieved IDs, name each local proxy explicitly, and leave source entailment and human usefulness to future held-out validation.

Evidence location within long contexts can materially affect model performance \citep{liu2024lostmiddle}. Source duplication can also distort training and evaluation \citep{lee2022dedup}. These findings motivate bounded context, documented truncation and source/video grouping in our pilot. Our lexical source-family graph is a conservative leakage guard, not proof that all source duplication or pretraining contamination has been removed.
